\section{Задание}

Разработать программное обеспечение с графическим интерфейсом для моделирования системы массового обслуживания (СМО) с использованием пошагового принципа ($\Delta t$) и событийного принципа. Программа должна определять максимальную длину очереди при заданных параметрах.

Рассматриваемая СМО состоит из генератора заявок, очереди ожидающих обработки заявок и обслуживающего аппарата (ОА). Программа должна предоставлять возможность выбора закона распределения для генератора и обслуживающего аппарата из следующих вариантов:

\begin{itemize}
    \item Равномерное распределение
    \item Экспоненциальное распределение
    \item Нормальное распределение
    \item Распределение Эрланга
\end{itemize}

Необходимо реализовать возможность ручного задания параметров распределений, количества заявок, временного шага и вероятности возврата обработанной заявки в очередь.

\section{Теоретическая часть}

\subsection{Принципы моделирования СМО}

\subsubsection{Пошаговый принцип ($\Delta t$)}

Принцип $\Delta t$ заключается в последовательном анализе состояний всех блоков системы в момент времени $t + \Delta t$ по заданному состоянию блоков в момент времени $t$. Новое состояние определяется в соответствии с алгоритмическим описанием блоков с учётом случайных факторов, задаваемых распределениями вероятности.

На каждом временном шаге проводится анализ, определяющий, какие события должны имитироваться в программной модели. Основной недостаток метода — значительные затраты машинного времени на анализ и контроль функционирования всей системы. При недостаточно малом $\Delta t$ возможен пропуск отдельных событий, что приводит к некорректным результатам моделирования.

\subsubsection{Событийный принцип}

Событийный принцип использует дискретный характер изменения состояний моделируемой системы. Состояния блоков имитационной модели анализируются только в момент появления события. Момент наступления следующего события определяется минимальным значением из списка будущих событий, представляющего совокупность ближайших изменений состояния каждого блока системы.

Этот подход более эффективен по сравнению с пошаговым, так как устраняет необходимость проверки состояния системы на каждом временном шаге.

\subsection{Используемые законы распределения}

\subsubsection{Равномерное распределение}

Случайная величина $X$ имеет равномерное распределение на отрезке $[a, b]$, если её плотность распределения $f(x)$ равна:

\begin{equation}
f(x) = \begin{cases}
\frac{1}{b-a}, & a \leq x \leq b \\
0, & \text{иначе}
\end{cases}
\end{equation}

Функция распределения $F(x)$ имеет вид:

\begin{equation}
F(x) = \begin{cases}
0, & x < a \\
\frac{x-a}{b-a}, & a \leq x \leq b \\
1, & x > b
\end{cases}
\end{equation}

Интервал времени между появлением $i$-й и $(i-1)$-й заявки по равномерному закону вычисляется следующим образом:

\begin{equation}
T_i = a + (b - a) \cdot R
\end{equation}

где $R$ — псевдослучайное число из интервала $[0, 1]$.

\subsubsection{Экспоненциальное распределение}

Экспоненциальное непрерывное распределение характеризуется плотностью:

\begin{equation}
f(x) = \lambda e^{-\lambda x}, \quad x \geq 0
\end{equation}

где $\lambda > 0$ — интенсивность (скорость наступления события).

Функция распределения:

\begin{equation}
F(x) = 1 - e^{-\lambda x}, \quad x \geq 0
\end{equation}

Интервал времени между появлением заявок по экспоненциальному закону вычисляется по формуле:

\begin{equation}
T_i = -\frac{1}{\lambda} \ln(1 - R)
\end{equation}

где $R$ — псевдослучайное число из интервала $[0, 1]$.

\subsubsection{Нормальное распределение}

Случайная величина $X$ имеет нормальное распределение с параметрами $\mu$ и $\sigma$, если её плотность распределения $f(x)$ равна:

\begin{equation}
f(x) = \frac{1}{\sigma\sqrt{2\pi}} e^{-\frac{(x-\mu)^2}{2\sigma^2}}, \quad x \in \mathbb{R}, \sigma > 0
\end{equation}

где $\mu$ — математическое ожидание (среднее значение), $\sigma$ — стандартное отклонение.

Функция распределения $F(x)$ выражается через функцию ошибок:

\begin{equation}
F(x) = \frac{1}{2}\left[1 + \text{erf}\left(\frac{x-\mu}{\sigma\sqrt{2}}\right)\right]
\end{equation}

Интервал времени между появлением заявок по нормальному закону вычисляется с использованием центральной предельной теоремы:

\begin{equation}
T_i = \sigma\sqrt{\frac{12}{n}}\left(\sum_{j=1}^{n}R_j - \frac{n}{2}\right) + \mu
\end{equation}

При $n=12$ формула упрощается:

\begin{equation}
T_i = \sigma\left(\sum_{j=1}^{12}R_j - 6\right) + \mu
\end{equation}

где $R_j$ — псевдослучайные числа из интервала $[0, 1]$.

\subsubsection{Распределение Эрланга}

Распределение Эрланга является частным случаем гамма-распределения с целым параметром формы. Плотность вероятности:

\begin{equation}
f(x) = \frac{\lambda^k x^{k-1} e^{-\lambda x}}{(k-1)!}, \quad x \geq 0
\end{equation}

где $k \in \mathbb{N}$ — параметр формы (число фаз), $\lambda > 0$ — интенсивность (скорость).

Функция распределения выражается через неполную гамма-функцию.

Интервал времени между появлением заявок по распределению Эрланга вычисляется как сумма $k$ экспоненциально распределённых величин:

\begin{equation}
T_i = -\frac{1}{k\lambda}\sum_{j=1}^{k}\ln(1 - R_j)
\end{equation}

где $R_j$ — псевдослучайные числа из интервала $[0, 1]$.

\section{Описание реализации}

\subsection{Используемые библиотеки}

Программа реализована на языке Python с использованием следующих библиотек:

\begin{itemize}
    \item \textbf{PyQt6} — создание графического интерфейса пользователя
    \item \textbf{random} — генерация псевдослучайных чисел
    \item \textbf{math} — математические функции (логарифм, квадратный корень)
\end{itemize}

\subsection{Структура программы}

Программа состоит из следующих модулей:

\begin{itemize}
    \item \texttt{main.py} — точка входа в приложение
    \item \texttt{gui/window.py} — реализация графического интерфейса
    \item \texttt{models/distributions.py} — классы для генерации случайных величин по различным распределениям
    \item \texttt{models/step\_model.py} — реализация пошагового подхода моделирования
    \item \texttt{models/event\_model.py} — реализация событийного подхода моделирования
    \item \texttt{constants.py} — константы интерфейса и параметры по умолчанию
\end{itemize}

\subsection{Метод решения}

Для каждого распределения реализован класс с методом \texttt{generate()}, возвращающим случайную величину в соответствии с формулами генерации, описанными в теоретической части.

Моделирование СМО выполняется двумя методами:

\textbf{Пошаговый метод:}
\begin{enumerate}
    \item Инициализация переменных: текущее время, время следующей генерации заявки, время окончания обработки, длина очереди
    \item Цикл моделирования с шагом $\Delta t$ до обработки заданного числа заявок
    \item На каждом шаге проверка событий генерации и завершения обработки
    \item Учёт вероятности возврата заявки в очередь
    \item Отслеживание максимальной длины очереди
\end{enumerate}

\textbf{Событийный метод:}
\begin{enumerate}
    \item Инициализация списка будущих событий с первым событием генерации
    \item Цикл моделирования до обработки заданного числа заявок:
    \begin{itemize}
        \item Выбор ближайшего события из списка
        \item Обработка события (генерация или окончание обработки)
        \item Добавление новых событий в список с сохранением сортировки по времени
        \item Учёт вероятности возврата заявки
        \item Отслеживание максимальной длины очереди
    \end{itemize}
\end{enumerate}

\section{Практическая часть}

На рисунке 1 представлен результат работы разработанной программы для СМО с выбранными параметрами:

\begin{itemize}
    \item Генератор: равномерное распределение с параметрами $a = 0$, $b = 10$
    \item Обслуживающий аппарат: нормальное распределение с параметрами $\mu = 5$, $\sigma = 2$
    \item Количество заявок: 1000
    \item Вероятность возврата: 10\%
    \item Временной шаг: 0.01
\end{itemize}

\includeimage
	{img1}{f}{H}{1.0\textwidth}
	{Результат работы программы}

Программа корректно вычисляет максимальную длину очереди для обоих методов моделирования. Различие в результатах пошагового и событийного подходов объясняется дискретизацией времени в пошаговом методе и зависит от величины временного шага $\Delta t$.

\section{Вывод}

В ходе выполнения лабораторной работы была разработана программа с графическим интерфейсом для моделирования системы массового обслуживания с использованием пошагового и событийного принципов.

Программа обеспечивает гибкий выбор законов распределения для генератора заявок и обслуживающего аппарата, автоматическую валидацию входных параметров и параллельное моделирование двумя методами с выводом результатов. Реализованный интерфейс позволяет проводить сравнительный анализ различных конфигураций СМО и исследовать влияние параметров распределений на характеристики системы.
